\documentclass[UTF8,a4paper,11pt]{ctexart}
\usepackage[margin=1.8cm]{geometry}
\usepackage{amsmath,amssymb,amsfonts,bm}
\usepackage{xcolor}
\usepackage{enumitem}
\usepackage{hyperref}
\hypersetup{colorlinks=true,linkcolor=blue,urlcolor=blue}

\newcommand{\bp}{\textcolor{red}{\textbf{【背】}}}
\newcommand{\uf}{\textcolor{blue}{\textbf{【了解】}}}
\newcommand{\km}{\textcolor{teal}{\textbf{【考点】}}}

\title{\textbf{理论力学复习提纲（精简版）}}
\author{燕山大学 · 应用物理学 · 理论力学}
\date{}

\begin{document}
\maketitle

\noindent\textbf{使用方法}：红字为“要背”，蓝字为“了解即可”。左边 $\square$ 用于背诵时打勾。对照详细版查推导。

% ===================== 一、必背清单 =====================
\section{必背清单（打勾复习）}

\subsection*{1. 约束与虚功原理 \km 概念1/4、选择1、大题5}
\begin{itemize}[label=$\square$]
  \item \bp 完整约束：可写成坐标有限形式；非完整约束：不可积分的速度约束（纯滚动、冰刀）。
  \item \bp 定常约束不显含 $t$；非定常显含 $t$。
  \item \bp 虚位移 $\delta\bm r_i$：给定瞬时、$\delta t=0$ 的约束允许无限小位移。
  \item \bp 理想约束：$\sum\bm N_i\cdot\delta\bm r_i=0$。
  \item \bp 虚功原理（平衡充要条件）：$\sum\bm F_i^{(a)}\cdot\delta\bm r_i=0$。
\end{itemize}

\subsection*{2. 拉格朗日方程 \km 选择2、大题3/4}
\begin{itemize}[label=$\square$]
  \item \bp $L=T-V$。
  \item \bp 拉氏方程：$\dfrac{d}{dt}\dfrac{\partial L}{\partial\dot q_j}-\dfrac{\partial L}{\partial q_j}=0$（保守力，$Q_j=0$）。
  \item \bp 广义动量 $p_j=\partial L/\partial\dot q_j$；循环坐标（缺 $q_k$）$\Rightarrow p_k=\text{const}$。
  \item \bp 解题四步：选坐标 $\rightarrow$ 写 $L$ $\rightarrow$ 代方程 $\rightarrow$ 化简。
\end{itemize}

\subsection*{3. 小振动与简正坐标 \km 概念3、大题1}
\begin{itemize}[label=$\square$]
  \item \bp 简正坐标：使 $T,V$ 同时对角化的坐标 $\xi$，每维对应一个独立简谐振动。
  \item \bp 耦合摆：$\xi_1=\theta_1+\theta_2,\ \xi_2=\theta_1-\theta_2$。
  \item \bp $\omega_1=\sqrt{g/l}$（同相），$\omega_2=\sqrt{g/l+2k/m}$（反相）。
\end{itemize}

\subsection*{4. 哈密顿力学 \km 概念2、选择3、大题3}
\begin{itemize}[label=$\square$]
  \item \bp $H=\sum p_j\dot q_j-L$（保守定常下 $H=T+V=E$）。
  \item \bp 正则方程：$\dot q_j=\partial H/\partial p_j,\ \dot p_j=-\partial H/\partial q_j$。
  \item \bp 相空间：以 $(q,p)$ 为坐标的 $2s$ 维空间，状态=相点。
  \item \bp 泊松括号：$\{u,v\}=\sum(\partial u/\partial q_j\,\partial v/\partial p_j-\partial u/\partial p_j\,\partial v/\partial q_j)$。
  \item \bp 正则括号：$\{q_j,p_k\}=\delta_{jk},\ \{q_j,q_k\}=0,\ \{p_j,p_k\}=0$。
  \item \bp 雅可比恒等式：$\{\{u,v\},w\}+\{\{v,w\},u\}+\{\{w,u\},v\}=0$。
\end{itemize}

\subsection*{5. 变分原理与正则变换 \km 选择3、大题2}
\begin{itemize}[label=$\square$]
  \item \bp 哈密顿原理：$\delta\!\int_{t_1}^{t_2}L\,dt=0$，端点 $\delta q(t_1)=\delta q(t_2)=0$。
  \item \bp 四类母函数：
    \begin{itemize}[label=$\square$]
      \item $F_1(q,Q,t)$：$p_j=\partial F_1/\partial q_j,\ P_j=-\partial F_1/\partial Q_j$
      \item $F_2(q,P,t)$：$p_j=\partial F_2/\partial q_j,\ Q_j=\partial F_2/\partial P_j$
      \item $F_3(p,Q,t)$：$q_j=-\partial F_3/\partial p_j,\ P_j=-\partial F_3/\partial Q_j$
      \item $F_4(p,P,t)$：$q_j=-\partial F_4/\partial p_j,\ Q_j=\partial F_4/\partial P_j$
    \item \textbf{记忆口诀}：自变量为 $(q,Q),(q,P),(p,Q),(p,P)$ 四组合；偏导符号固定——含$q$得$+p$、含$Q$得$-P$、含$p$得$-q$、含$P$得$+Q$。
    \end{itemize}
  \item \bp 哈密顿—雅可比方程：$H(q,\partial S/\partial q,t)+\partial S/\partial t=0$。
\end{itemize}

% ===================== 二、公式速查 =====================
\section{核心公式速查}
\[
\begin{array}{ll}
\hline
\text{名称} & \text{公式} \\
\hline
\text{虚功原理} & \sum\bm F_i^{(a)}\!\cdot\!\delta\bm r_i=0 \\[4pt]
\text{拉氏方程} & \dfrac{d}{dt}\dfrac{\partial L}{\partial\dot q_j}-\dfrac{\partial L}{\partial q_j}=0 \\[8pt]
\text{广义动量} & p_j=\dfrac{\partial L}{\partial\dot q_j} \\[6pt]
\text{哈密顿量} & H=\sum p_j\dot q_j-L \\[6pt]
\text{正则方程} & \dot q_j=\dfrac{\partial H}{\partial p_j},\ \dot p_j=-\dfrac{\partial H}{\partial q_j} \\[8pt]
\text{泊松括号} & \{u,v\}=\sum\!\left(\dfrac{\partial u}{\partial q_j}\dfrac{\partial v}{\partial p_j}
-\dfrac{\partial u}{\partial p_j}\dfrac{\partial v}{\partial q_j}\right) \\[10pt]
\text{哈密顿原理} & \delta\!\int_{t_1}^{t_2}L\,dt=0 \\[8pt]
\text{母函数(新H)} & K=H+\dfrac{\partial F}{\partial t} \\[6pt]
\text{H-J方程} & H\!\left(q,\dfrac{\partial S}{\partial q},t\right)+\dfrac{\partial S}{\partial t}=0 \\[10pt]
\text{阿特伍德机} & \ddot x=\dfrac{M-m}{M+m}g \\[8pt]
\text{耦合摆频率} & \omega_1=\sqrt{g/l},\ \omega_2=\sqrt{g/l+2k/m} \\[6pt]
\hline
\end{array}
\]

% ===================== 三、了解即可 =====================
\section{了解即可（不要求死背推导）}
\begin{itemize}[itemsep=2pt]
  \item \uf 达朗贝尔原理：$\sum(\bm F_i^{(a)}-m_i\bm a_i)\cdot\delta\bm r_i=0$（拉氏方程出发点）。
  \item \uf 广义能量积分：$h=\sum\dot q_j\partial L/\partial\dot q_j-L=T+V=E$（定常+保守）。
  \item \uf 刘维尔定理：相体积守恒，$d\rho/dt=0$。
  \item \uf 位力定理：$2\langle T\rangle=n\langle V\rangle$（$V\propto r^n$）。
  \item \uf 正则变换：保持正则方程形式的坐标变换。
  \item \uf 欧拉—拉格朗日方程：$\partial F/\partial y-d/dx(\partial F/\partial y')=0$。
\end{itemize}

% ===================== 四、题型对策 =====================
\section{题型对策与易错}
\begin{enumerate}[itemsep=3pt]
  \item \km \textbf{概念/选择}：完整vs非完整（能否写坐标有限式）、相空间、简正坐标、虚功原理、哈密顿原理变量 $q,p,t$。
  \item \km \textbf{选择2陷阱}：拉氏方程“不显含某 $q_k$”（不是不显含 $\dot q_k$）$\Rightarrow$ 广义动量守恒。
  \item \km \textbf{大题1 耦合摆}：取 $\xi_1=\theta_1+\theta_2,\ \xi_2=\theta_1-\theta_2$ 解耦，求两频率。
  \item \km \textbf{大题2 母函数}：写出母函数定义 $\rightarrow$ 证变换后正则（$\{Q,Q\}=0,\{P,P\}=0,\{Q,P\}=\delta$ 或 $K=H+\partial F/\partial t$）。
  \item \km \textbf{大题3 滑轮三法}：牛顿（受力图）/拉氏（$L=T-V$）/哈密顿（$p\to H\to$正则方程），结果一致 $\ddot x=(M-m)g/(M+m)$。
  \item \km \textbf{大题4 两自由度拉氏}：思路简单（照四步），难在代数，耐心展开。
  \item \km \textbf{大题5 虚功}：选坐标 $\rightarrow$ 写虚位移 $\rightarrow$ 代虚功原理 $\rightarrow$ 系数归零。
  \item \textbf{母函数正负号}：$p_j=+\partial F/\partial q_j,\ P_j=-\partial F/\partial Q_j$（第一类）；$F_2$ 用 $Q_j=+\partial F_2/\partial P_j$。
\end{enumerate}

\end{document}
